\input amstex 
\input opmac

% Document setup
\fontfam[LM fonts]
\typosize[13/16]
\frenchspacing
\parskip10pt
\parindent0pt
\abovedisplayskip12pt
\abovedisplayshortskip0pt
\belowdisplayskip12pt
\belowdisplayshortskip12pt
\medskipamount10pt
\raggedbottom

% AmsTeX environments
\TagsOnRight

% OpMaC hyperlink colors
\hyperlinks{\Green}{\Blue}

% Toggle to display solutions of examples / proofs of theorems
\newif\ifDisplayTexts
\DisplayTextstrue

% Prints the title of the handout
\def\Title#1{\centerline{\typosize[18/]\bf #1}\medskip}
% Prints the author of the handout
\def\Author#1{\centerline{\typosize[14/]{\em #1}}\bigskip}

% Prints a link to the MathComps website to the handout
\def\MathcompsLink#1{%
    \centerline{\url{https://mathcomps.fun/\link/#1}}%
    \vskip2pt
    \centerline{\typosize[12/]{\em(\linkCaption)}}%
    \vskip10pt
}%

% Helper macro to create a link using OpMaC
\def\Link[#1]#2{\ulink[#1]{#2}}

% Defines the Czech translations of the captions
\def\definecaptionsCS{%
	\cslang%
    \def\captionExercise{Cvičení}%
    \def\captionTheorem{Tvrzení}%
    \def\captionProof{Důkaz}%
    \def\captionDefinition{Definice}%
    \def\captionExample{Příklad}%
    \def\captionSolution{Řešení}%
    \def\captionProblem{Úloha}%
    \def\captionProblemSolutions{Řešení k úlohám}
    \def\captionExerciseSolutions{Řešení ke cvičením}
    \def\captionHints##1{%
      \ifcase##1
        % Case 0: TeX starts counting at 0
      \or
        První%
      \or 
        Druhé%
      \or 
        Třetí%
      \or 
        Čtvrté%
      \or 
        Páté%
      \or 
        Šesté%
      \else
        \errmessage{Too many hints lol}%
      \fi
      \ návody k řešením%
    }
    \def\link{cs/materialy}
    \def\linkCaption{interaktivní online verze}
}

% Defines the Slovak translations of the captions
\def\definecaptionsSK{%
	\sklang%
    \def\captionExercise{Cvičenie}%
    \def\captionTheorem{Tvrdenie}%
    \def\captionProof{Dôkaz}%
    \def\captionDefinition{Definícia}%
    \def\captionExample{Príklad}%
    \def\captionSolution{Riešenie}%
    \def\captionProblem{Úloha}%
    \def\captionProblemSolutions{Riešenia k úlohám}
    \def\captionExerciseSolutions{Riešenia k cvičeniam}
    \def\captionHints##1{%
      \ifcase##1
        % Case 0: TeX starts counting at 0
      \or
        Prvé%
      \or 
        Druhé%
      \or 
        Tretie%
      \or 
        Štvrté%
      \or 
        Piate%
      \or 
        Šieste%
      \else
        \errmessage{Too many hints lol}%
      \fi
      \space návody k riešeniam%
    }
    \def\link{sk/materialy}
    \def\linkCaption{interaktívna online verzia}
}

% Defines the English translations of the captions
\def\definecaptionsEN{%
	\uslang%
    \def\captionExercise{Exercise}%
    \def\captionTheorem{Theorem}%
    \def\captionProof{Proof}%
    \def\captionDefinition{Definition}%
    \def\captionExample{Example}%
    \def\captionSolution{Solution}%
    \def\captionProblem{Problem}%
    \def\captionProblemSolutions{Solutions to problems}
    \def\captionExerciseSolutions{Solutions to exercises}
    \def\captionHints##1{%
      \ifcase##1
        % Case 0: TeX starts counting at 0
      \or
        First%
      \or 
        Second%
      \or 
        Third%
      \or 
        Fourth%
      \or 
        Fifth%
      \or 
        Sixth%
      \else
        \errmessage{Too many hints lol}%
      \fi
      \space hints%
    }
    \def\link{en/handouts}
    \def\linkCaption{interactive online version}
}

% Sets the language of the handout
\def\setlanguage#1{%
    \csname definecaptions#1\endcsname
}

% The default language is Slovak
\setlanguage{SK}

% Convenience macro for starting an item inside OPMac's itemize
\let\i\startitem

% Redefine parentheses, math convenience
\def\({\left(}
\def\){\right)}

% Helper to print the text on the very right side of the page
\def\flushright#1{{\unskip\nobreak\hskip\parfillskip\penalty50
  \hskip0.5em\hbox{}\nobreak\hfil #1
  \parfillskip=0pt \finalhyphendemerits=0 \par}}

% Prints the qedsquare symbol
\def\qedsquare{\flushright{\vrule height 5pt width 5pt depth 0pt\relax}}

% Quick mathbb symbols
\def\defbb #1{\expandafter\long\expandafter\def\csname#1\endcsname{{\bbchar #1}}}
\defbb N
\defbb Z 
\defbb Q 
\defbb C
\defbb R

% Quick bold, italic, etc. commands, LaTeX style
\def\textbf#1{\begingroup\bf #1\endgroup}
\def\textit#1{\begingroup\it #1\endgroup}
\def\textrm#1{\begingroup\rm #1\endgroup}
\def\texttt#1{\begingroup\tt #1\endgroup}
\def\textsf#1{\begingroup\sf #1\endgroup}
\def\textsl#1{\begingroup\sl #1\endgroup}

% Helper macros to get rid of spaces when reading/printing
\long\def\saferead#1{\ignorespaces#1}
\long\def\safepar{\par\removelastskip}

% Prints the end of a block (exercise, theorem, example, definition)
\def\endblock{%
  \safepar\medskip%
  \centerline{%
    \hbox{%
      \vrule width 2em height 0.4pt depth 0pt
      \hskip 0.75em
      $\diamond$%
      \hskip 0.75em
      \vrule width 2em height 0.4pt depth 0pt
    }%
  }%
  \safepar\medskip%
}

% Prints a left-aligned bar with the given text
\long\def\leftbar#1{% 
  \safepar\medskip\medskip%
  \hbox{%
    \vrule width 2pt
    \hskip 0.75em
    \vtop{%
      \hsize=\dimexpr\hsize-1.5em-2pt\relax
      \parindent=1em
      \noindent #1\par
    }%
  }%
  \safepar\medskip\medskip%
}

% Counters for already definied exercises, theorems etc
\newcount\exercises
\newcount\theorems
\newcount\problems
\newcount\examples
\newcount\definition

% #1 - Source (might be empty)
% #2 - Statement
% #3 - Solution
\long\def\Exercise#1#2#3{%
	\advance\exercises by 1\relax%
	\expandafter\long\expandafter\def\csname e\the\exercises-solution\endcsname{#3}%
	\leftbar{%
		{\bf \captionExercise\ \the\exercises}%
		\def\temp{#1}\ifx\temp\empty\else\ (\temp)\fi.%
		\enskip\saferead{#2}\safepar%
	}%
}

% #1 - Source (might be empty)
% #2 - Statement
% #3 - Solution
\long\def\Theorem#1#2#3{%
	\advance\theorems by 1\relax%
	\leftbar{%
		{\bf \captionTheorem\ \the\theorems}%
		\def\temp{#1}\ifx\temp\empty\else\ (\temp)\fi.%
		\enskip\saferead{#2}\safepar%
	}%
	\safepar{\em \captionProof.}\enskip\saferead{#3}\qedsquare%
	\endblock%
}

% #1 - Source (might be empty)
% #2 - Statement
% #3 - Solution
\long\def\Example#1#2#3{%
	\advance\examples by 1\relax%
	\leftbar{%
		{\bf \captionExample\ \the\examples}%
		\def\temp{#1}\ifx\temp\empty\else\ (\temp)\fi.%
		\enskip\saferead{#2}%
	}%
	\ifDisplayTexts%
		\safepar{\em \captionSolution.}\enskip\saferead{#3}\endblock%
	\else%
		\safepar\medskip%	
	\fi%
}

% A helper macro which prints the given number of stars
\newcount\starcount
\def\Stars#1{%
  \ifnum#1>0%
    \def\starholder{}%
    \starcount=#1%
    \loop\ifnum\starcount>0
      \expandafter\def\expandafter\starholder\expandafter{\starholder\star}%
      \advance\starcount by-1
    \repeat
    ${}^{\starholder}$%
  \fi
}

% The maximal number of hints some problem has
% This is used to determine how many hint sections we need
\newcount\maxHints
\maxHints=0

% Helper to update max hints
\def\UpdateMaxHints#1{%
  \ifnum#1>\maxHints%
    \maxHints=#1%
  \fi%
}

% The local variable used to count the args after statement
% These are the hints and the solution
\newcount\problemArgsCount

% #1 - Number Of Stars
% #2 - Source (might be empty)
% #3 - Statement
% Then we have an arbitrary number of hints
% And then the solution
\long\def\Problem#1#2#3{%
  \advance\problems by 1\relax%
  \expandafter\def\csname p\the\problems-difficulty\endcsname{#1}%
  \problemArgsCount=0\relax%
  \leftbar{%
    {\bf \captionProblem\ \the\problems\Stars{#1}}%
    \def\temp{#2}\ifx\temp\empty\else\ (\temp)\fi.%
    \enskip\saferead{#3}\safepar
  }%
  \ProblemReadArgs%
}


% Helper to identify spaces
\def\problemSpace{ }

% The macro which starts the search for args of the \Problem macro
\long\def\ProblemReadArgs{%
  \futurelet\next\ProblemCheckNext%
}

% Check what the next token is in the current reading of the \Problem
\def\ProblemCheckNext{%
  \ifx\next\bgroup%
    % If it is a brace {, read the argument
    \expandafter\ProblemReadOneArg%
  \else%
    % If it is a space, eat it and try again
    \ifx\next\problemSpace%
      \expandafter\expandafter\expandafter\ProblemIgnoreSpace%
    \else%
      % Stop reading, we hit something that isn't { or space.
      % Now distribute the generic 'args' into 'hints' and 'solution'
      \ProblemDistributeArgs%
    \fi%
  \fi%
}

% A helper macro which just eats a space and continues with reading problem args
\def\ProblemIgnoreSpace#1{%
  \ProblemReadArgs%
}

% Reads the argument and continues searching for more
\long\def\ProblemReadOneArg#1{%
  \advance\problemArgsCount by 1\relax%
  \expandafter\long\expandafter\def\csname p\the\problems-arg\the\problemArgsCount\endcsname{#1}%
  \ProblemReadArgs%
}

% A helper index
\newcount\currentArgIndex

% Remembers hints and solution of the current problem
\def\ProblemDistributeArgs{%
  % Handle no args...At least one should be there, the solution
  \ifnum\problemArgsCount=0%
    \errmessage{A problem with no solution?}
  % We have some args
  \else%
    % Set the solution (the last arg)
    \expandafter\let\expandafter\tempContent\csname p\the\problems-arg\the\problemArgsCount\endcsname%
    \expandafter\let\csname p\the\problems-solution\endcsname\tempContent%
    %
    % Set the hints
    % Loop from 1 to (Count - 1)
    \currentArgIndex=1\relax%
    \loop\ifnum\currentArgIndex<\problemArgsCount%
      % Get content of pX-argY
      \expandafter\let\expandafter\tempContent\csname p\the\problems-arg\the\currentArgIndex\endcsname%
      % Define pX-hintY with that content
      \expandafter\let\csname p\the\problems-hint\the\currentArgIndex\endcsname\tempContent%
      % Move on to the next arg
      \advance\currentArgIndex by 1\relax%
    \repeat%
    %
    % Update max hints
    \UpdateMaxHints{\the\numexpr\problemArgsCount-1\relax}%
  \fi%
}

% Counter for currently displayed solutions
\newcount\displayedProblemSolutions

% Displays the section with all problem solutions
\def\DisplayProblemSolutions{%
  % Use Opmac \sec for the header
  \csname\string\sec:M\endcsname{\captionProblemSolutions}
  %
  % Display each solution
	\loop\ifnum\displayedProblemSolutions<\problems%
		\advance\displayedProblemSolutions by 1\relax%
		{\bf \captionProblem\ \the\displayedProblemSolutions.} %
		\csname p\the\displayedProblemSolutions-solution\endcsname\endgraf%
	\repeat%	
}

% The counter for displayed exercise solutions 
\newcount\displayedExerciseSolutions

% Displays the section with all exercise solutions
\def\DisplayExerciseSolutions{%
  % Use Opmac \sec for the header
  \csname\string\sec:M\endcsname{\captionExerciseSolutions}
  % 
  % Display each solution
	\loop\ifnum\displayedExerciseSolutions<\exercises%
		\advance\displayedExerciseSolutions by 1\relax%
		{\bf \captionExercise\ \the\displayedExerciseSolutions.}\enskip%
		\csname e\the\displayedExerciseSolutions-solution\endcsname\endgraf%
	\repeat%	
}


% The current displayed level of hints
\newcount\hintLevel

% Entry point for displaying all hints recursively
\def\DisplayHints{%
  \DisplayHintsLoop%
}

% Recursive Loop
\def\DisplayHintsLoop{%
  % If we haven't displayed all hint levels yet
  \ifnum\hintLevel<\maxHints%
    % Increment which level we're displaying
    \advance\hintLevel by 1\relax%
    %
    % Use Opmac \sec for the header
    \csname\string\sec:M\endcsname{\captionHints{\the\hintLevel}}%
    %
    % Call the inner loop (which processes all problems for this level)
    \DisplayHintsLevel{\the\hintLevel}%
    %
    % Recursive call
    \DisplayHintsLoop%
  \fi%
}

% The current problem whose hints are being displayed
\newcount\currentProblemNumber%

% Display all hints of a specific level
\def\DisplayHintsLevel#1{%
  % Reset current problem for this level
  \currentProblemNumber=0%
  % Display all hints
  \loop\ifnum\currentProblemNumber<\problems%
    \advance\currentProblemNumber by 1\relax%
    % Check if this specific hint exists
    \expandafter\ifx\csname p\the\currentProblemNumber-hint#1\endcsname\relax%
      % Do nothing if hint doesn't exist for this problem
    \else%
      {\bf \captionProblem\ \the\currentProblemNumber.}\enskip%
      \csname p\the\currentProblemNumber-hint#1\endcsname\safepar%
    \fi%
  \repeat%
}

% Main image macro: \Image{path}{scale} or \Image{path} (defaults to scale=1)
\def\Image#1{%
  \isnextchar\bgroup{\ImageWithScale{#1}}{\ImageWithScale{#1}{1}}%
}

% Internal macro that does the actual work
\def\ImageWithScale#1#2{%
  % #1 = filename path
  % #2 = scale factor
  %
  \edef\imagepath{Images/#1}%
  % First, load the PDF image to measure its natural dimensions
  \pdfximage{\imagepath}%
  %
  % Place the image in box 0 at natural size to get its dimensions
  \setbox0=\hbox{\pdfrefximage\pdflastximage}%
  %
  % Extract natural width
  \dimen0=\wd0
  %
  % Extract natural height (height + depth for total vertical size)
  \dimen2=\ht0
  \advance\dimen2 by \dp0
  %
  % Scale the width by the scale factor
  \dimen0=#2\dimen0
  %
  % Scale the height by the scale factor
  \dimen2=#2\dimen2
  %
  % Load the image again with scaled dimensions
  \pdfximage width \dimen0 height \dimen2 {\imagepath}%
  %
  % Center the image horizontally using \line with \hss
  \safepar\medskip%
  \line{\hss\pdfrefximage\pdflastximage\hss}%
  \safepar\medskip%
}

% Highlight macro: makes a paragraph stand out with a box
\long\def\Highlight#1{%
  % Create some vertical space before the highlight
  \safepar\medskip\vskip\parskip%
  %
  % Start a group to keep settings local
  \bgroup%
  %
  % Set up the frame parameters
  \setbox0=\vbox{%
    % Add horizontal padding inside the box
    \hsize=\dimexpr\hsize-2em\relax%
    % Add some vertical padding
    \kern0.5em%
    % Indent the text slightly
    \leftskip=1em%
    \rightskip=1em%
	% This is to negate the parskip which will come
	\vskip-\parskip%
    % The actual content
    #1\safepar%
    % Bottom padding
    \kern0.5em%
  }%
  %
  % Center the box horizontally and add a frame around it
  \line{\hss%
    \vbox{%
      % Draw a frame around box0
      \hrule%
      \hbox{%
        \vrule%
        \vbox{%
          \box0%
        }%
        \vrule%
      }%
      \hrule%
    }%
  \hss}%
  %
  \egroup%
  %
  % Add some vertical space after the highlight
  \safepar\medskip
}